\section{Conclusion}

We asked whether register tokens regularize a vision transformer when data is
scarce, and answered it with a fully paired 24-run controlled ablation across two
distinct data regimes (100 and 300 images per class) in which the register count
was the strictly isolated variable. The answer is no. No register configuration
improves held-out Top-1 accuracy over the register-free baseline ($75.19 \pm 0.24\%$
at 100 images/class and $81.33 \pm 0.37\%$ at 300 images/class). While tripling the
training budget provides a clear $+6.14\%$ accuracy gain, registers confer no
advantage over control in either regime ($p > 0.50$), confirming that the empirical
null result is scale-invariant across low-data regimes.

The study's more transferable methodological result is how nearly an uncareful
protocol would have concluded the opposite. On the 100-image/class validation split,
$K=4$ narrows the generalization gap in every seed and lowers validation loss in
every seed at $p = 0.024$, presenting what initially appears to be a textbook
regularization signature. Crucially, this effect completely vanishes on the held-out
test split ($p = 0.823$). The training loss, invariant across all four arms to within
$0.005$ nats, reveals the mechanism: nothing was traded during optimization, so
nothing was regularized. A generalization gap measured on the very split used for
model selection does not constitute evidence of generalization.

At $K=8$, we find the single consistent held-out effect in the low-data regime, and
it is a penalty: test loss degrades across all three seeds ($p = 0.014$) and
final-block attention entropy declines in all three. On a 192-dimensional backbone,
allocating eight tokens that carry no visual content imposes a measurable capacity cost.

We view this as a scope result rather than a refutation of the register mechanism.
Registers were introduced as a pretraining-time intervention at foundation scale;
we applied them at fine-tuning time to a compact pretrained model under constrained
data budgets. Mechanistic attention diagnostics confirm that pretrained routing does
not reorganize to exploit registers in low-data regimes. The definitive experiment to
distinguish whether registers do not help in small models or whether they cannot be
added post-hoc is a from-scratch replication at this scale, and our open sweep
infrastructure executes that evaluation without modification.
